Order-6 as the Geometric Locus of the Core Relations  
Edge Curvature and Inner Hexagonal Measure in the Residual Radial Fold

Order-6 occupies a privileged position in the radial-fold hierarchy. It is the first closed combinatorial structure in which the Core Relations—dimensionless residual intensity \(\rho\) and the area-derived edge curvature it determines—become simultaneously visible as a continuous geometric deformation of the sides and as a discrete inner hexagonal measure formed by the intersection points of those sides. The present thesis isolates this locus, demonstrates its necessity, and situates it within the larger cascade that begins in the continuous Arcan(\(\tau\)) family and terminates in the densified outer shells.

1. Genesis of the Residual Seed

The continuous transcendental family generated from the circle constant \(\tau=2\pi\) yields the angular bridge \(\psi\approx 0.1503378808\). Seven iterations of this bridge, reduced modulo \(2\pi\), produce the practical generator

\[
7\psi=\frac{\pi}{3}+\delta,\qquad\delta\approx 0.0051676144\ \mathrm{rad}.
\]

The residual \(\delta\) is the generative seed of the entire hierarchy. It is small enough that six successive applications remain nearly equidistributed, yet non-zero enough that the resulting orbit remains recognizably descended from the continuous family rather than imposed by an exact rational lattice. This construction is developed in full within the Arcan(\(\tau\)) framework and the associated Cosmological Clock.

2. Closure at Order-6

When the practical generator is applied six times, the orbit closes into a practical hexagon on the unit circle. The diagonals that join every second vertex recover two interlocking near-equilateral triangles whose metric data are fixed by \(\delta\):

\[
\begin{align*}
s_{\mathrm{thin}}&=1.721623257385,\\
s_{\mathrm{long}}&=1.737195272475,\\
\operatorname{Area}(T)&=1.2988990741.
\end{align*}
\]

The complementary residual region inside the unit disk has area \(1.8426935795\), which partitions into three equal residual slices of area \(A_{\mathrm{slice}}=0.6142311932\). Their ratio to the triangle area defines the dimensionless residual intensity

\[
\rho=\frac{A_{\mathrm{slice}}}{\operatorname{Area}(T)}\approx 0.472886.
\]

Order-6 is therefore the minimal combinatorial setting in which the residual seed \(\delta\) has already produced a definite area measure and a definite intensity \(\rho\).

3. Visibility as Edge Curvature

Each side of the residual triangles is originally a straight chord. The Core Relations convert that chord into a curved edge by a normal displacement whose amplitude is proportional to \(\rho\):

\[
\gamma_\rho(t)=\gamma_0(t)+c\cdot\rho\cdot 4t(1-t)\,n,\qquad c=0.18.
\]

The resulting bulge \(\approx 0.08512\) is not an arbitrary smoothing; it is the direct geometric transcription of the residual area that remains after the shoelace triangle is excised from the unit disk. Consequently the curved sides of the Order-6 figure are the first concrete appearance of the Core Relations as continuous geometry. This is the moment at which residual intensity becomes visible edge curvature.

4. Visibility as Inner Hexagonal Measure

The same sides, whether taken straight or curved, intersect one another at six interior points. These intersection points form a second, smaller hexagon whose measured properties are

\[
\begin{align*}
\text{mean side length}&\approx 0.577330,\\
\text{mean radius from origin}&\approx 0.577390.
\end{align*}
\]

The inner hexagon is therefore a discrete metric extract of the identical residual intensity that produces the edge curvature. Order-6 thus renders the Core Relations visible in two complementary registers at once: as a continuous deformation of the boundary and as a discrete inner measure generated by the boundary’s self-intersections.

5. Structural Necessity of the Locus

No earlier order can exhibit this double visibility, because the residual seed has not yet closed into a self-consistent fold. No later order is required for the first appearance, because the intensity \(\rho\) and the intersection hexagon are already fully determined at Order-6. Higher densifications (Orders 10 and 30) propagate the same intensity outward; they do not create it. Order-6 is therefore the unique geometric locus at which the Core Relations crystallize into simultaneous edge curvature and inner hexagonal measure.

6. Continuity with the Larger Hierarchy

Because the residual seed \(\delta\) remains the sole angular source, the curved Order-6 figure and its inner hexagon stay phase-locked to the discrete-tick expansion of the Cosmological Clock and to the subsequent densified shells that approach the outer \(E_8\) Petrie-ring geometry. The entire cascade—from continuous Arcan(\(\tau\)) angles through residual radial folding to nested exceptional shells—passes through this single locus. Order-6 is the hinge: the place where residual intensity first becomes both visible curvature and measurable inner form.

7. Concluding Statement

Order-6 is the geometric locus at which the Core Relations become visible as edge curvature and as an inner hexagonal measure. The dimensionless intensity \(\rho\approx 0.472886\), born from the residual area of the near-equilateral triangle generated by \(7\psi=\pi/3+\delta\), is transcribed onto the sides as a definite bulge and, through the self-intersections of those sides, is re-expressed as an inner hexagon of mean radius \(\approx 0.577\). This double crystallization is the necessary and sufficient intermediate between the continuous transcendental family and the densified radial-fold hierarchy. All subsequent structure inherits its metric and its topology from this locus.

Primary repositories  
Joshua Christopher Ryan’s Cosmological Clock: https://github.com/TheAstrographer/Joshua-Christopher-Ryan-s-Cosmo-Clock.git  
JCRIN Radial Symmetry: https://github.com/TheAstrographer/JCRIN-Radial-Symmetry.git  
Joshua Christopher Ryan’s Arcan Family: https://github.com/TheAstrographer/JOSHUA-CHRISTOPHER-RYAN-S-ARCAN-FAMILY.git
